\documentclass[10pt,a4paper]{amsart}
\usepackage[margin=19mm]{geometry}
\usepackage{amsmath,amssymb,mathtools}
\usepackage[scr=rsfs,cal=euler]{mathalfa}
\usepackage{xfrac}
\usepackage[unicode,hidelinks,bookmarks=false]{hyperref}
\newcommand{\ie}{i.e.\ }
\newcommand{\cf}{cf.\ }
\newcommand{\resp}{respectively\ }
\newcommand{\Subsection}{\subsection}
\providecommand{\nobreakdash}{\nobreak\hbox{-}}
\setlength{\emergencystretch}{2em}
\multlinegap0pt
\newcommand\version{July 28, 2026}
\usepackage{lmodern}
\usepackage{mathtools}
\usepackage{enumitem}
\usepackage{cleveref}
\newtheorem{theorem}{Theorem}[section]
\newtheorem{proposition}[theorem]{Proposition}
\newtheorem{lemma}[theorem]{Lemma}
\newtheorem{corollary}[theorem]{Corollary}
\newtheorem{claim}[theorem]{Claim}
\newtheorem{definition}[theorem]{Definition}
\newtheorem{problem}[theorem]{Problem}
\newtheorem{remark}[theorem]{Remark}
\newcommand{\T}{\mathbb T}
\newcommand{\Sone}{\mathbb S^1}
\newcommand{\Z}{\mathbb Z}
\newcommand{\R}{\mathbb R}
\newcommand{\N}{\mathbb N}
\newcommand{\C}{\mathbb C}
\newcommand{\dd}{\,\mathrm d}
\newcommand{\avgint}{\mathop{\rlap{\raisebox{.13ex}{--}}\!\int}}
\newcommand{\wh}{\widehat}
\newcommand{\seminorm}[1]{\left[#1\right]}
\DeclareMathOperator{\degmap}{deg}
\DeclareMathOperator{\supp}{supp}
\hypersetup{
  pdftitle={No Universal Fourier Summation Formula for the Degree of Circle Maps},
  pdfauthor={Rupert L. Frank and Paata Ivanisvili},
  pdfsubject={Fourier summation processes and the winding number},
  pdfkeywords={winding number, topological degree, Fourier coefficients, summation process, fractional Sobolev space, Hölder space}
}
\begin{document}
\title[The Winding Number at the Critical H"older Exponent 1/3: Fai]{The Winding Number at the Critical H\"older Exponent 1/3: Failure of Universal Fourier Summation}
\author{Munich Center for Quantum Science and Technology, Schellingstr. 4, 80799 München, Germany; Secondary 42A16, 46E35, 55M25 ęywords Winding number, topological degree, Fourier coefficients, summation process, fractional Sobolev space, Hölder space, autocorrelation e̱gin document; Rupert L. Frank, Paata Ivanisvili}
\date{}
\maketitle
\noindent\emph{Source and licence.} The Winding Number at the Critical H\"older Exponent 1/3: Failure of Universal Fourier Summation. Version arXiv:2607.25862v1 (CC BY 4.0). This material is distributed under the Creative Commons Attribution 4.0 International licence, http://creativecommons.org/licenses/by/4.0/, as recorded in the arXiv OAI licence metadata for this exact version and shown on the abstract page https://arxiv.org/abs/2607.25862v1. Full rights notice: licenses/arxiv-2607.25862-CC-BY-4.0.txt.\par\smallskip
\noindent\textit{Editorial scope.} Counted unit: the original \texttt{theorem} environment \texttt{thm:synthesis} at source lines 235--246 together with its complete proof at lines 654--716; the delivered document keeps source lines 227--907 so that every referenced label and every citation resolves inside it. Native editable source is the paper's own concatenated TeX; the AMS wrapper is editorial formatting only. No publisher class, upstream script or unrelated artwork is redistributed.
\section{Synthesis by autocorrelations of small phases}\label{sec:synthesis}

The purpose of this section is to prove the following synthesis theorem for odd periodic functions with zero derivative at the origin in terms of autocorrelation functions. The imaginary part of the autocorrelation function of $f\in L^2(\T)$ is denoted by
\begin{equation}
    \label{eq:imag-autocorrelation}
    B_f(t) :=\operatorname{Im}\int_\T f(x+t)\overline{f(x)}\,\frac{dx}{2\pi} \qquad\text{for}\ t\in\T \,.
\end{equation}


\begin{theorem}[Small-phase autocorrelation synthesis]\label{thm:synthesis}
Let $P$ be a smooth real odd $2\pi$-periodic function with $P'(0)=0$.  Given $\delta>0$, there are real coefficients $(\Lambda_\ell)_{\ell\geq1}$ and smooth real periodic functions $(\theta_\ell)_{\ell\geq1}$ such that
\begin{align}
 &\sum_{\ell\geq1}|\Lambda_\ell|<\infty,\label{eq:synthesis-l1}\\
 &\|\theta_\ell\|_{C^1}<\delta
 \quad\text{for every }\ell,\label{eq:synthesis-small}\\
 &P(t)=\sum_{\ell\geq1}\Lambda_\ell
 B_{e^{i\theta_\ell}}(t)
 \quad\text{for every }t\in\T.\label{eq:synthesis-identity}
\end{align}
The series in \eqref{eq:synthesis-identity} converges absolutely and uniformly.
\end{theorem}


The proof shows that the smoothness assumption on $P$ can be relaxed to $P\in Y_0$, where $Y_0$ is introduced in the next subsection. We will only apply this theorem in the smooth situation.

This whole section is dedicated to the proof of this theorem.

%%%%%%%%%%%%%%%%%%%%%%%%%%%%%%%%%%%%%%

\subsection{A weighted sine space}

Let $Y$ be the Banach space of real odd functions
\begin{equation}\label{eq:Y-definition}
 H(t)=\sum_{n\geq1}h_n\sin(nt)
 \quad\text{such that}\quad
 \|H\|_Y:=\sum_{n\geq1}n^3|h_n|<\infty.
\end{equation}
The series and its first derivative converge absolutely and uniformly.  Define the closed subspace
\begin{equation}\label{eq:Y0-definition}
 Y_0:=\{H\in Y:H'(0)=0\}
 =\left\{H\in Y:\sum_{n\geq1}nh_n=0
 \right\}.
\end{equation}
For positive integers $r,s$, set
\begin{equation}\label{eq:Ers}
 E_{r,s}(t):=\sin(rt)+\sin(st)-\sin((r+s)t)
 \qquad\text{for all}\ t\in\T \,.
\end{equation}
We note that $E_{r,s}\in Y_0$.

\begin{lemma}[Balanced atomic decomposition]\label{lem:balanced-decomp}
There are a countable index set $\Gamma$ and, for each $\gamma\in\Gamma$, positive integers $r_\gamma$ and $s_\gamma$ satisfying
\[
 r_\gamma\leq s_\gamma\leq r_\gamma+1 \,,
\]
and a linear map $Y_0\ni H\mapsto (a_\gamma(H))_{\gamma\in\Gamma}$ such that, for all $H\in Y_0$,
\begin{equation}\label{eq:atomic-decomposition}
 H=\sum_{\gamma\in\Gamma} a_\gamma(H)E_{r_\gamma,s_\gamma}
\end{equation}
with convergence in $Y$ and, with $N_\gamma:=r_\gamma+s_\gamma$, 
\begin{equation}\label{eq:atomic-bound}
 \sum_{\gamma\in\Gamma} |a_\gamma(H)|N_\gamma^3
 \leq\frac32\|H\|_Y.
\end{equation}
\end{lemma}

In what follows we will call a pair $(r,s)$ of positive integers \emph{balanced} if $r\leq s\leq r+1$. For example, the sum in \eqref{eq:atomic-decomposition} is over balanced pairs.

\begin{proof}
\emph{Step 1.} We expand $H$ in its sine series, $H(t)=\sum_{n\geq1}h_n\sin(nt)$. Since $H\in Y_0$, we have
\[
 h_1=-\sum_{n\geq2}nh_n \,.
\]
This allows us to write
\begin{equation}\label{eq:H-by-G}
 H=\sum_{n\geq2}h_nG_n \,,
\end{equation}
where, for $n\geq1$, we have set
\[
 G_n(t):=\sin(nt)-n\sin t \,.
\]
Note in particular, that $G_1=0$. In view of \eqref{eq:H-by-G} it suffices to decompose each $G_n$ into atoms $E_{r,s}$ with balanced pairs $(r,s)$.

\medskip

\emph{Step 2.} We shall show that for every $n$ there is a finite index set $\mathcal T_n$ such that
\begin{equation}\label{eq:G-tree}
    G_n=-\sum_{v\in\mathcal T_n}E_{r_v,s_v} \,.
\end{equation}

The basic ingredient in the proof of \eqref{eq:G-tree} is the elementary identity
\begin{equation}\label{eq:G-recursion}
 G_n=G_r+G_s-E_{r,s}
 \qquad\text{where}\ n= r+s \,.
\end{equation}

Before proving \eqref{eq:G-tree} for general $n$, we treat some basic examples. For $n=2$, we have, using $G_1=0$,
$$
G_2 = G_1 + G_1 - E_{1,1} = - E_{1,1} \,.
$$
Thus, $\mathcal T_2$ consists of a single element and $(r_v,s_v)=(1,1)$ for this element $v$. Next, for $n=3$, we have, using the formulas for $G_1$ and $G_2$, 
$$
G_3 = G_1 + G_2 - E_{1,2} = - E_{1,1} - E_{1,2} \,.
$$
Thus, $\mathcal T_3$ consists of two elements, and $(r_v,s_v)$ takes the two values $(1,1)$ and $(1,2)$. Next, for $n=4$, we have, using the formula for $G_2$,
$$
G_4 = G_2 + G_2 - E_{2,2} = - E_{1,1} - E_{1,1} - E_{2,2} \,.
$$
We emphasize that $\mathcal T_4$ consists of \emph{three} elements and that the pair $(1,1)$ appears \emph{twice}, while $(2,2)$ appears once.

We now generalize this procedure to arbitrary integers $n\geq 2$. We recursively split every integer $N\geq2$ into
\[
 r=\lfloor N/2\rfloor,
 \qquad
 s=\lceil N/2\rceil,
\]
and stop at the leaves labeled $1$.  Denote by $\mathcal T_n$ the set of internal vertices of this finite binary tree, with the two child subtrees regarded as disjoint copies when they have the same label.  Iterating \eqref{eq:G-recursion} and using $G_1=0$ yields the claimed decomposition \eqref{eq:G-tree}.

\medskip

\emph{Step 3.}
Denoting by $N_v=r_v+s_v$ the label at an internal vertex $v\in\mathcal T_n$, we claim that
\begin{equation}\label{eq:tree-cost}
 C_n:=\sum_{v\in\mathcal T_n}N_v^3 \leq\frac32n^3 \,.
\end{equation}
The assertion is trivial for $n=1$.  For $n\geq2$, with the balanced split $n=r+s$,
\[
 C_n=n^3+C_r+C_s.
\]
By induction,
\[
 C_n\leq n^3+\frac32(r^3+s^3).
\]
For a balanced pair, $r^3+s^3\leq n^3/3$.  When $n$ is even this is immediate; when $n\geq3$ is odd,
\[
 r^3+s^3=\frac{n^3+3n}{4}\leq\frac{n^3}{3}.
\]
This proves \eqref{eq:tree-cost}.

\medskip

\emph{Step 4.}
In Steps 1 and 2 we have shown that $H$ has a decomposition of the form \eqref{eq:atomic-decomposition} with 
$$
\Gamma:=\{ (n,v) :\ n\geq 2 \,,\ v\in\mathcal T_n \}
$$
and with coefficients $a_{(n,v)}(H)= - h_n$ depending linearly on $H$.

Now we use the bound from Step 3 to prove \eqref{eq:atomic-bound} and convergence in $Y$. Indeed, \eqref{eq:tree-cost} gives
\[
 \sum_\gamma|a_\gamma(H)|N_\gamma^3
 \leq\sum_{n\geq2}|h_n|C_n
 \leq\frac32\sum_{n\geq2}n^3|h_n| \leq \frac32 \|H\|_Y \,.
\]
Finally,
\begin{equation}
    \label{eq:normers}
    \|E_{r,s}\|_Y=r^3+s^3+(r+s)^3\leq2(r+s)^3,
\end{equation}
with the same bound when $r=s$ after combining the two equal sine terms.  Hence the expansion converges absolutely in $Y$.
\end{proof}

Different indices $\gamma$ may correspond to the same atom $E_{r_\gamma,s_\gamma}$, so representations in this redundantly indexed family are not intrinsically unique. For instance, the two distinct child vertices in the tree for $G_4$ both correspond to $E_{1,1}$. The construction above nevertheless fixes a particular bounded linear coefficient-selection map, satisfying \eqref{eq:atomic-bound}, and this fixed choice is used below.

%%%%%%%%%%%%%%%%%%%%%%%%%%%%%%%%%%

\subsection{Exact cubic atoms}

For every balanced pair $(r,s)$, we define
\begin{equation}\label{eq:phi-rs}
 \phi_{r,s}(x):=
 \begin{cases}
 \displaystyle
 \frac{\cos(rx)}{r}+\frac{\cos(sx)}{s}
 +\frac{\sin((r+s)x)}{r+s},&r<s,\\[2.2ex]
 \displaystyle
 \frac{\cos(rx)}{r}+\frac{\sin(2rx)}{2r},&r=s.
 \end{cases}
\end{equation}

\begin{lemma}[Cubic autocorrelation atoms]\label{lem:cubic-atoms}
For every balanced pair $(r,s)$,
\begin{equation}\label{eq:cubic-atom}
 \int_\T
 \big(\phi_{r,s}(x+t)-\phi_{r,s}(x)\big)^3\,\frac{dx}{2\pi}
 =\lambda_{r,s}E_{r,s}(t),
\end{equation}
where
\begin{equation}\label{eq:lambda-rs}
 \lambda_{r,s}:=
 \begin{cases}
 \displaystyle-\frac{3}{rs(r+s)},&r<s,\\[1.5ex]
 \displaystyle-\frac{3}{4r^3},&r=s.
 \end{cases}
\end{equation}
Moreover, for $0\leq k\leq3$, we have, setting $N:=r+s$,
\begin{equation}\label{eq:phi-derivative-bound}
 \|\phi_{r,s}^{(k)}\|_\infty\leq C N^{k-1}.
\end{equation}
In particular,
\begin{equation}\label{eq:lambda-comparable}
 cN^{-3}\leq|\lambda_{r,s}|\leq CN^{-3}.
\end{equation}
\end{lemma}

\begin{proof}
The derivative bounds follow immediately from \eqref{eq:phi-rs} and the balance relations $r,s\simeq N$.

We prove \eqref{eq:cubic-atom} by Fourier expansion.  If
\[
 \phi(x)=\sum_{k\in\Z}\phi_k e^{ikx},
\]
then
\begin{equation}\label{eq:cubic-fourier-general}
 \int_\T(\phi(x+t)-\phi(x))^3\,\frac{dx}{2\pi}
 =\sum_{k_1+k_2+k_3=0}
 \phi_{k_1}\phi_{k_2}\phi_{k_3}
 \prod_{j=1}^3(e^{ik_jt}-1).
\end{equation}
For positive $r,s$ and $N=r+s$, one has
\begin{equation}\label{eq:triangle-product}
 (e^{irt}-1)(e^{ist}-1)(e^{-iNt}-1)
 =2iE_{r,s}(t).
\end{equation}

Suppose first that $r<s$.  The relevant Fourier coefficients are
\[
 \phi_{\pm r}=\frac1{2r},
 \qquad
 \phi_{\pm s}=\frac1{2s},
 \qquad
 \phi_N=-\frac{i}{2N},
 \qquad
 \phi_{-N}=\frac{i}{2N}.
\]
The only nonvanishing resonant contribution to \eqref{eq:cubic-fourier-general} comes from permutations of $(r,s,-N)$ and $(-r,-s,N)$.  There is one apparent exception: when $(r,s)=(1,2)$, the triples $(r,r,-s)$ and $(-r,-r,s)$ are also resonant.  Their common Fourier coefficient is real, while
\[
 (e^{irt}-1)^2(e^{-i2rt}-1)=2iE_{r,r}(t),
 \qquad
 (e^{-irt}-1)^2(e^{i2rt}-1)=-2iE_{r,r}(t).
\]
Thus these two contributions cancel exactly.  Using the six permutations of each of the two main resonances and \eqref{eq:triangle-product}, we obtain
\begin{align*}
 &6\left(\frac1{2r}\right)
 \left(\frac1{2s}\right)
 \left(\frac{i}{2N}\right)(2iE_{r,s})\\
 &\qquad+
 6\left(\frac1{2r}\right)
 \left(\frac1{2s}\right)
 \left(-\frac{i}{2N}\right)(-2iE_{r,s})
 =-\frac{3}{rsN}E_{r,s}.
\end{align*}

If $r=s$, the only resonance is $(r,r,-2r)$ and its conjugate.  There are three permutations of each.  Since
\[
 \phi_{\pm r}=\frac1{2r},
 \qquad
 \phi_{2r}=-\frac{i}{4r},
 \qquad
 \phi_{-2r}=\frac{i}{4r},
\]
formula \eqref{eq:triangle-product} gives
\[
 -\frac{3}{4r^3}E_{r,r}(t).
\]
This proves \eqref{eq:cubic-atom}.  The comparability \eqref{eq:lambda-comparable} follows from balance.
\end{proof}

%%%%%%%%%%%%%%%%%%%%%%%%%%%%%%%%%%%%%%%%%%%%

\subsection{A uniform fifth-order estimate}

For a balanced pair $(r,s)$, we set
$$
b_{r,s}(\tau,t) := B_{e^{i\tau\phi_{r,s}}}(t) = \operatorname{Im}\int_\T
 e^{i\tau(\phi_{r,s}(x+t)-\phi_{r,s}(x))}\,\frac{dx}{2\pi}.
$$

\begin{lemma}[Fifth-order bound]\label{lem:fifth-order}
There exists a universal constant $C$ such that for every balanced pair $(r,s)$ and every $|\tau|\leq1$, we have, with $N:=r+s$,
\begin{equation}\label{eq:fifth-order-bound}
 \left\|\frac{\partial^5}{\partial\tau^5}
 b_{r,s}(\tau)\right\|_Y
 \leq C N^{-2}.
\end{equation}
\end{lemma}

In the proof we will use the following observation concerning the imaginary part of the autocorrelation function. By Parseval's identity, we have, for each $h\in L^2(\T)$,
    $$
    \int_\T h(x+t) \overline{h(x)}\,\frac{dx}{2\pi} = \sum_{n\in\Z} |\wh h(n)|^2 e^{int} \,.
    $$
    The series on the right side is absolutely and uniformly convergent. Taking the imaginary part, we find
    \begin{equation}\label{eq:B-sine-series}
        B_h(t) = \sum_{n\in\Z} |\wh h(n)|^2 \sin(nt) = \sum_{n\geq 1} \left( |\wh h(n)|^2 - |\wh h(-n)|^2 \right) \sin(nt) \,.
    \end{equation}


\begin{proof}
We abbreviate $\phi:=\phi_{r,s}$ and $f_\tau:=e^{i\tau\phi}$.  Put
\[
 u_a:=\phi^a e^{i\tau\phi},
 \qquad 0\leq a\leq5.
\]
We first justify differentiation in the weighted Fourier space used below. For
smooth functions $u,v$, Cauchy--Schwarz gives
\begin{equation}\label{eq:weighted-bilinear}
 \sum_{n\in\Z}|n|^3
 |\wh u(n)\overline{\wh v(n)}|
 \leq \|u\|_{\dot H^3}\|v\|_2.
\end{equation}
Moreover, for every $0\leq a\leq5$, the map
$\tau\mapsto \phi^a e^{i\tau\phi}$ is smooth with values in both $H^3$ and
$L^2$. Thus \eqref{eq:weighted-bilinear} and the Banach-valued product rule
show that the Fourier-energy sequence
$\tau\mapsto (|\wh{f_\tau}(n)|^2)_{n\in\Z}$ is five times continuously
differentiable with values in
$\ell^1(\Z\setminus\{0\};|n|^3)$. Differentiation gives
\begin{equation}\label{eq:fifth-energy-derivative}
    \frac{d^5}{d\tau^5} \, |\wh{f_\tau}(n)|^2
 =\sum_{a=0}^5\binom5a i^a(-i)^{5-a}
 \wh u_a(n)\,\overline{\wh u_{5-a}(n)} \,.
\end{equation}
It follows from \eqref{eq:B-sine-series} with $h=f_\tau$ that the sine coefficient of $b_{r,s}(\tau)$ at frequency $n\geq1$ is $|\wh f_\tau(n)|^2 - |\wh f_\tau(-n)|^2$.  Hence
\begin{align}
 \left\|\frac{\partial^5}{\partial\tau^5}
 b_{r,s}(\tau)\right\|_Y
 &\leq\sum_{n\in\Z}|n|^3 \left|\frac{d^5}{d\tau^5} \, |\wh{f_\tau}(n)|^2 \right| \notag\\
 &\lesssim \sum_{a=0}^5
 \left(\sum_n|n|^6|\wh u_a(n)|^2\right)^{1/2}
 \left(\sum_n|\wh u_{5-a}(n)|^2\right)^{1/2}\notag\\
 &=\sum_{a=0}^5
 \|u_a\|_{\dot H^3}\,\|u_{5-a}\|_2 \,.
 \label{eq:Y-by-Sobolev}
\end{align}
We used Cauchy--Schwarz in the second line.

The derivative bounds \eqref{eq:phi-derivative-bound}, the product rule, and $|\tau|\leq1$ imply
\begin{equation}\label{eq:ua-estimates}
 \|u_a\|_{H^3}\lesssim N^{3-a},
 \qquad
 \|u_a\|_2\lesssim N^{-a},
 \qquad \text{for all}\ 0\leq a\leq5.
\end{equation}
For completeness, we prove these bounds.  For $0\leq j\leq3$, every term obtained by applying $j$ derivatives to $\phi^a e^{i\tau\phi}$ is a constant multiple of
\[
 \phi^{a-q}
 \prod_{\nu=1}^{q}\phi^{(k_\nu)}
 \prod_{\mu=1}^{v}\phi^{(\ell_\mu)}
 e^{i\tau\phi},
\]
where $0\leq q\leq a$, all $k_\nu,\ell_\mu\geq1$, and
\[
 \sum_{\nu=1}^{q}k_\nu+
 \sum_{\mu=1}^{v}\ell_\mu=j.
\]
The first product records derivatives falling on $\phi^a$, while the second records those generated by differentiating the exponential.  By \eqref{eq:phi-derivative-bound}, its absolute value is at most a constant times
\[
 N^{-(a-q)}
 N^{\sum_\nu(k_\nu-1)}
 N^{\sum_\mu(\ell_\mu-1)}
 = N^{j-a-v} \leq N^{j-a} \,.
\]
Taking $j=0,1,2,3$ proves \eqref{eq:ua-estimates}.  

Thus every product in \eqref{eq:Y-by-Sobolev} is bounded by a constant times 
\[
 N^{3-a}N^{-(5-a)}= N^{-2} \,.
\]
This proves \eqref{eq:fifth-order-bound}.
\end{proof}

\begin{lemma}[Approximate triangle atoms]\label{lem:approx-atoms}
For $0<\varepsilon\leq1/2$, define
\begin{equation}\label{eq:Etilde}
 \widetilde E_{r,s}^{\varepsilon}
 :=\frac{2B_{e^{i\varepsilon\phi_{r,s}}}
 -B_{e^{2i\varepsilon\phi_{r,s}}}}
 {\varepsilon^3\lambda_{r,s}}.
\end{equation}
Then $\widetilde E_{r,s}^{\varepsilon}\in Y_0$ and
\begin{equation}\label{eq:Etilde-error}
 \|\widetilde E_{r,s}^{\varepsilon}-E_{r,s}\|_Y
 \leq C\varepsilon^2(r+s).
\end{equation}
\end{lemma}

\begin{proof}
The map $\tau\mapsto b_{r,s}(\tau)$ is odd because
\[
 B_{e^{-i\tau\phi}}=-B_{e^{i\tau\phi}}.
\]
Consequently,
\[
 b_{r,s}(0)=b_{r,s}^{(2)}(0)=b_{r,s}^{(4)}(0)=0.
\]
Furthermore,
\[
 b_{r,s}'(0,t)
 =\int_\T(\phi(x+t)-\phi(x))\,\frac{dx}{2\pi}=0,
\]
and Lemma~\ref{lem:cubic-atoms} gives
\[
 b_{r,s}^{(3)}(0)=-\lambda_{r,s}E_{r,s}.
\]
The Banach-valued Taylor formula and Lemma~\ref{lem:fifth-order} therefore yield
\begin{equation}\label{eq:b-taylor}
 \left\|b_{r,s}(\tau)
 +\frac{\tau^3\lambda_{r,s}}6E_{r,s}\right\|_Y
 \leq C|\tau|^5N^{-2}
 \qquad(|\tau|\leq1).
\end{equation}
The cubic terms in $2b(\varepsilon)-b(2\varepsilon)$ combine to $\varepsilon^3\lambda_{r,s}E_{r,s}$.  Applying \eqref{eq:b-taylor} at $\varepsilon$ and $2\varepsilon$, dividing by $\varepsilon^3|\lambda_{r,s}|$, and using \eqref{eq:lambda-comparable} gives \eqref{eq:Etilde-error}.

Finally, if $f=e^{i\theta}$ with a smooth periodic real phase function, then
\[
 B_f'(0)=\operatorname{Im}\int_\T f'(x)\overline{f(x)}\,\frac{dx}{2\pi}
 =\int_\T\theta'(x)\,\frac{dx}{2\pi}=0.
\]
Moreover, $f$ and $B_f$ are smooth, so the sine coefficients of $B_f$ have
the weighted summability required in the definition of $Y$. Thus each term
in \eqref{eq:Etilde} lies in $Y_0$.
\end{proof}

%%%%%%%%%%%%%%%%%%%%%%%%%%%%%%%%%%

\subsection{Proof of the small-phase autocorrelation synthesis theorem}

We are now in a position to prove the main result of this section.


\begin{proof}[Proof of Theorem \ref{thm:synthesis}]
\emph{Step 1.}
According to Lemma~\ref{lem:balanced-decomp}, we have a bounded linear coefficient-selection map $Y_0\ni H \mapsto (a_\gamma(H))$. In terms of this map, for $0<\varepsilon\leq1/2$ and $H\in Y_0$, set
\begin{equation}\label{eq:S-epsilon}
 S_\varepsilon H
 :=\sum_{\gamma\in\Gamma} a_\gamma(H)
 \widetilde E_{r_\gamma,s_\gamma}^{\varepsilon}.
\end{equation}
This series converges absolutely in $Y$. Indeed, by \eqref{eq:normers} and \eqref{eq:Etilde-error},
\[
 \|\widetilde E_{r,s}^{\varepsilon}\|_Y
 \leq \|E_{r,s}\|_Y+C\varepsilon^2(r+s)
 \leq C(r+s)^3,
\]
and therefore \eqref{eq:atomic-bound} implies
\[
 \sum_\gamma |a_\gamma(H)|
 \|\widetilde E_{r_\gamma,s_\gamma}^{\varepsilon}\|_Y
 <\infty.
\]
By Lemma \ref{lem:approx-atoms}, we have $\widetilde E_{r,s}^{\varepsilon}\in Y_0$, so the closedness of $Y_0$ implies that $S_\varepsilon H\in Y_0$. To summarize, $S_\varepsilon$ is a bounded linear operator from $Y_0$ to $Y_0$.

By \eqref{eq:Etilde-error} and \eqref{eq:atomic-bound},
\begin{align*}
 \|(S_\varepsilon-I)H\|_Y
 &\leq C\varepsilon^2
 \sum_\gamma|a_\gamma(H)|N_\gamma\\
 &\leq C\varepsilon^2
 \sum_\gamma|a_\gamma(H)|N_\gamma^3
 \leq C\varepsilon^2\|H\|_Y.
\end{align*}
Choose $0<\varepsilon\leq1/2$ so small that
\begin{equation}\label{eq:Neumann-small}
 \|S_\varepsilon-I\|_{Y_0\to Y_0} \leq \frac12
\end{equation}
and, in addition, $2C_0\varepsilon<\delta$, where $C_0:= \sup_{(r,s) \,\mathrm{balanced}} \|\phi_{r,s}\|_{C^1}$, which is finite by Lemma~\ref{lem:cubic-atoms}.  Then $S_\varepsilon$ is invertible by the Neumann series.  

\medskip

\emph{Step 2.}
Let $P$ be a smooth, real, odd, $2\pi$-periodic function with $P'(0)=0$. Then $P\in Y_0$. Put
\[
 H:=S_\varepsilon^{-1}P \,.
\]
By \eqref{eq:S-epsilon} and \eqref{eq:Etilde},
\begin{equation}\label{eq:P-expanded-B}
 P
 =\sum_\gamma
 \frac{a_\gamma(H)}{\varepsilon^3\lambda_{r_\gamma,s_\gamma}}
 \left(2B_{e^{i\varepsilon\phi_{r_\gamma,s_\gamma}}}
 -B_{e^{2i\varepsilon\phi_{r_\gamma,s_\gamma}}}
 \right).
\end{equation}
The scalar coefficients in this expansion are summable.  Indeed, by \eqref{eq:lambda-comparable} and \eqref{eq:atomic-bound},
\begin{equation}\label{eq:coefficient-total-variation}
 \sum_\gamma
 \frac{3|a_\gamma(H)|}{\varepsilon^3|\lambda_{r_\gamma,s_\gamma}|}
 \leq C\varepsilon^{-3}
 \sum_\gamma|a_\gamma(H)|N_\gamma^3
 <\infty.
\end{equation}
Since $|B_f(t)|\leq1$, expansion \eqref{eq:P-expanded-B} therefore converges absolutely and uniformly after its two terms are enumerated as a single sequence.  Regrouping the two consecutive terms associated with each atom recovers the $Y$-convergent grouped series defining $S_\varepsilon H=P$. Since the split series converges unconditionally and uniformly, its uniform sum is therefore exactly $P$.  Finally, the phases in \eqref{eq:P-expanded-B} are $\varepsilon\phi_{r,s}$ and $2\varepsilon\phi_{r,s}$; their $C^1$ norms are smaller than $\delta$ by our choice of $\varepsilon$.
\end{proof}


%%%%%%%%%%%%%%%%%%%%%%%%%%%%%%%%%%%%%%%%%%%%%
%%%%%%%%%%%%%%%%%%%%%%%%%%%%%%%%%%%%%%%%%%%%%

\section{Kahane's cubic function and the critical barycenter}\label{sec:kahane}

The purpose of this section is to prove the following proposition. We recall that the notation $B_f$ was introduced in \eqref{eq:imag-autocorrelation}.

\begin{proposition}[Critical barycenter]\label{prop:critical-barycenter}
For every $r>0$, there are real numbers $(\Lambda_\ell)_{\ell\geq1}$ with
\begin{equation}\label{eq:Lambda-l1}
 \sum_{\ell\geq1}|\Lambda_\ell|<\infty \,,
\end{equation}
a family of real functions
\[
 \eta_{\ell,\tau}\in C^{0,1/3}(\T),
 \qquad \ell\geq1,\ \tau\in\T \,,
\]
such that, for every $\ell$, the map $(\tau,x)\mapsto\eta_{\ell,\tau}(x)$ is continuous, with
\begin{equation}\label{eq:small-barycenter-phases}
 \|\eta_{\ell,\tau}\|_{C^{0,1/3}}<r,
\end{equation}
and a constant $\kappa\ne0$ such that, for $h_{\ell,\tau}:=e^{i\eta_{\ell,\tau}}$, one has $\degmap h_{\ell,\tau}=0$ and
\begin{equation}\label{eq:critical-barycenter}
 \sum_{\ell\geq1}\Lambda_\ell
 \int_\T B_{h_{\ell,\tau}}(t)\,\frac{d\tau}{2\pi}
 =\kappa\sin t
 \qquad\text{for every }t\in\T.
\end{equation}
\end{proposition}

To prove this proposition, we will apply Theorem~\ref{thm:synthesis} with a particular function $P$. That function will be constructed in terms of a function that Kahane used in a related context \cite{Kahane2005}.

Before presenting the details of this argument, we deduce a simple consequence of it, which will be crucial in the proof of our main result.

For a bounded sequence $q:\Z\to\C$ and $f\in L^2(\T)$, set
\begin{equation}\label{eq:Qq}
 Q_q(f):=\sum_{n\in\Z}q(n)|\wh f(n)|^2.
\end{equation}
This series is absolutely convergent. If $q$ is real-valued and odd, then
$Q_q(f)$ is real and
\begin{equation}\label{eq:Qq-odd}
 Q_q(f)
 =\sum_{n\geq1}q(n)
 \big(|\wh f(n)|^2-|\wh f(-n)|^2\big).
\end{equation}

\begin{corollary}\label{cor:barycenter-multiplier}
Given $r>0$, let $(\Lambda_\ell)_{\ell\geq1}$, $(h_{\ell,\tau})_{\ell\geq 1, \tau\in\T}$ and $\kappa\neq 0$ be as in Proposition~\ref{prop:critical-barycenter}. Then every bounded real odd sequence $s$ satisfies
\begin{equation}\label{eq:barycenter-Q}
 \sum_{\ell\geq1}\Lambda_\ell
 \int_\T Q_s(h_{\ell,\tau})\,\frac{d\tau}{2\pi}
 =\kappa s(1).
\end{equation}
\end{corollary}

\begin{proof}
    Equation \eqref{eq:imag-autocorrelation} implies that, for each $h\in L^2(\T)$,
    $$
    2 \int_\T B_h(t) \sin(nt)\,\frac{dt}{2\pi} = |\wh h(n)|^2 - |\wh h(-n)|^2 
    \qquad\text{for all}\, n\geq1 \,.
    $$

    We apply this identity to $h=h_{\ell,\tau}$, integrate with respect to $\tau$, multiply by $\Lambda_\ell$ and sum over $\ell$ to obtain
    $$
    2 \int_\T \left( \sum_{\ell} \Lambda_\ell \int_\T B_{h_{\ell,\tau}}(t)\,\frac{d\tau}{2\pi} \right) \sin(nt)\,\frac{dt}{2\pi} = \sum_{\ell} \Lambda_\ell \int_\T \left( |\wh h_{\ell,\tau}(n)|^2 - |\wh h_{\ell,\tau}(-n)|^2 \right)\frac{d\tau}{2\pi}
    \quad\text{for all}\ n\geq 1\,.
    $$
    On the left side we interchanged the $t$-integration with the $\tau$-integration and the $\ell$-summation, which is justified since the series in \eqref{eq:critical-barycenter} converges uniformly because of $|B_{h_{\ell,\tau}}|\leq1$ and \eqref{eq:Lambda-l1}.

    Proposition \ref{prop:critical-barycenter} allows us to rewrite this as
    $$
    \kappa \mathbf 1_{\{n=1\}} = 2\kappa \int_\T \sin(t) \sin(nt)\,\frac{dt}{2\pi}
    = \sum_{\ell} \Lambda_\ell \int_\T \left( |\wh h_{\ell,\tau}(n)|^2 - |\wh h_{\ell,\tau}(-n)|^2 \right)\frac{d\tau}{2\pi}
    \qquad\text{for all}\ n\geq 1    \,.
    $$
    Multiplying by $s(n)$ and summing over $n\geq1$ gives \eqref{eq:barycenter-Q}. Here the interchange of the $n$-summation and the $\ell$-summation and the $\tau$-integration is justified by boundedness of $s$, \eqref{eq:Lambda-l1}, and $\sum_{n\in\Z}|\wh h_{\ell,\tau}(n)|^2=1$.
\end{proof}





%%%%%%%%%%%%%%%%%%%%%%%%%%%%

\subsection{The critical cubic identity}

Let
\begin{equation}\label{eq:kahane-function}
 \psi(x):= \frac{2^{1/9}}{3^{1/3}} \,\sum_{j=0}^{\infty}2^{-j/3}\sin(2^j x) \qquad\text{for}\ x\in\T \,.
\end{equation}

The following lemma is due to Kahane
\cite{Kahane2005}. We use the corrected lower summation
index $j=0$: in \cite[Eq.~(13)]{Kahane2005} it is printed as $j=1$, while
\cite[Eq.~(17)]{Kahane2005} and the direct computation below show that
$j=0$ is required. We include a proof for completeness.

\begin{lemma}\label{lem:kahane}
The series \eqref{eq:kahane-function} defines a function in $C^{0,1/3}(\T;\R)$ and
\begin{equation}\label{eq:kahane-identity}
 \int_\T
 \big(\psi(x+t)-\psi(x)\big)^3\,\frac{dx}{2\pi}
 =\sin t
 \qquad\text{for all}\ t\in\T \,.
\end{equation}
\end{lemma}

\begin{proof}
Uniform convergence is immediate from $\sum_j2^{-j/3}<\infty$.  To prove the H\"older bound, fix $0<|h|\leq1$ and choose $m\geq0$ so that
\[
 2^{-m-1}<|h|\leq2^{-m}.
\]
For the low frequencies, we bound
\begin{align*}
 \sum_{j=0}^m2^{-j/3}
 |\sin(2^j(x+h))-\sin(2^j x)|
 &\leq |h|\sum_{j=0}^m2^{2j/3}\\
 &\lesssim |h|2^{2m/3}
 \lesssim |h|^{1/3}.
\end{align*}
For the high frequencies, we bound
\[
 \sum_{j>m}2^{-j/3}
 |\sin(2^j(x+h))-\sin(2^j x)|
 \leq2\sum_{j>m}2^{-j/3}
 \lesssim 2^{-m/3}
 \lesssim |h|^{1/3}.
\]
This proves the required estimate for $0<|h|\leq1$. For increments with
$|h|>1$, we have
\[
 |\psi(x+h)-\psi(x)|
 \leq2\|\psi\|_\infty
 \leq2\|\psi\|_\infty |h|^{1/3},
\]
so the estimate holds for all increments on $\T$. Thus
$\psi\in C^{0,1/3}$.

Let $A:= \frac{2^{1/9}}{3^{1/3}}$ and
\[
 \psi_L(x):=A\sum_{j=0}^L2^{-j/3}\sin(2^j x) \,.
\]
In the integral of $(\psi_L(x+t)-\psi_L(x))^3$, a triple of dyadic frequencies can sum to zero only through
\[
 2^j+2^j=2^{j+1},
 \qquad 0\leq j<L.
\]
Indeed, if one power of two is the sum of two others, uniqueness of binary expansion forces the two smaller powers to be equal.  We now compute the contribution of this resonance.  Put $k:=2^j$.  The elementary formula
\[
 \sin(k(x+t))-\sin(kx)
 =2\sin(kt/2)\cos(kx+kt/2)
\]
and the identity
\[
 \int_\T\cos^2(kx+a)\cos(2kx+2a)\,\frac{dx}{2\pi}=\frac14
\]
give
\begin{align*}
 &\int_\T
 \big(\sin(k(x+t))-\sin(kx)\big)^2
 \big(\sin(2k(x+t))-\sin(2kx)\big)\,\frac{dx}{2\pi}\\
 &\qquad=2\sin^2(kt/2)\sin(kt)
 =\sin(kt)-\frac12\sin(2kt).
\end{align*}
The multinomial coefficient is $3$, and the product of the three lacunary amplitudes is
\[
 A^3 2^{-2j/3}2^{-(j+1)/3}=A^3 2^{-j-1/3}.
\]
Summing over $j$ therefore gives
\begin{align}
 &\int_\T
 \big(\psi_L(x+t)-\psi_L(x)\big)^3\,\frac{dx}{2\pi}\notag\\
 &\qquad=\frac32A^3 2^{-1/3}
 \sum_{j=0}^{L-1}2^{-j}
 \big(2\sin(2^j t)-\sin(2^{j+1}t)\big).\label{eq:finite-kahane}
\end{align}
The sum telescopes:
\[
 \sum_{j=0}^{L-1}2^{-j}
 \big(2\sin(2^j t)-\sin(2^{j+1}t)\big)
 =2\big(\sin t-2^{-L}\sin(2^L t)\big).
\]
Since $3A^3 2^{-1/3}=1$, we obtain the exact finite identity
\begin{equation}\label{eq:finite-kahane-exact}
 \int_\T
 \big(\psi_L(x+t)-\psi_L(x)\big)^3\,\frac{dx}{2\pi}
 =\sin t-2^{-L}\sin(2^L t).
\end{equation}
The functions $\psi_L$ converge uniformly to $\psi$, so the cubic integrands converge uniformly in $x$ for each $t$; in fact the convergence is uniform in $(x,t)$.  Passing to the limit in \eqref{eq:finite-kahane-exact} proves \eqref{eq:kahane-identity}.
\end{proof}

\begin{thebibliography}{99}
\bibitem[Kahane2005]{Kahane2005} Jean-Pierre Kahane, \emph{Sur l'\'{e}quation fonctionnelle $\int_{\mathbb T}(\psi(t+s)-\psi(s))^3\,ds=\sin t$}, C. R. Math. Acad. Sci. Paris \textbf{341} (2005), no.~3, 141--145, \href{https://doi.org/10.1016/j.crma.2005.05.016}{doi:10.1016/j.crma.2005.05.016}.
\end{thebibliography}
\end{document}
